% --------------------------------------------------------------------
%  LISTING STYLES
% --------------------------------------------------------------------
%  Two base looks, and one named style per language built on them, so
%  the colours and frame are the same in every book:
%
%      utpcode          the base for source code
%      utpterminal      the base for commands typed and output printed
%      utppython        Python
%      utpr             R
%      utpc             C
%      utpcpp           C++
%      utpbash          shell scripts and commands
%      utpcnumbered     C with line numbers, for a complete source file
%      utpbashnumbered  shell with line numbers, likewise
%
%      \begin{lstlisting}[style=utppython] ... \end{lstlisting}
%      \lstinputlisting[style=utpcnumbered]{all_experiments/Lab04/fcfs.c}
%
%  A book that sets most of its listings in one style makes it the
%  default once, after the input:  \lstset{style=utpcpp}

% Resolve repeated colour mixtures once, preserving the original RGB values.
\colorlet{UTPListingString}{UTPGreen!70!black}
\colorlet{UTPListingBackground}{UTPPaleGray!72}
\colorlet{UTPListingCodeFrame}{UTPBlue!28}
\colorlet{UTPListingTerminalFrame}{UTPMuted!35}

\lstdefinestyle{utpcode}{
  basicstyle=\ttfamily\scriptsize,
  keywordstyle=\bfseries\color{UTPBlue},
  commentstyle=\itshape\color{UTPMuted},
  stringstyle=\color{UTPListingString},
  backgroundcolor=\color{UTPListingBackground},
  frame=single,
  rulecolor=\color{UTPListingCodeFrame},
  framesep=5pt,
  breaklines=true,
  breakatwhitespace=false,
  showstringspaces=false,
  keepspaces=true,
  columns=fullflexible,
  upquote=true,
  tabsize=4,
  aboveskip=0.7em,
  belowskip=0.8em
}
\lstdefinestyle{utpterminal}{
  basicstyle=\ttfamily\scriptsize,
  backgroundcolor=\color{UTPListingBackground},
  frame=single,
  rulecolor=\color{UTPListingTerminalFrame},
  framesep=5pt,
  breaklines=true,
  showstringspaces=false,
  keepspaces=true,
  columns=fullflexible,
  upquote=true,
  aboveskip=0.6em,
  belowskip=0.7em
}

\lstdefinestyle{utppython}{style=utpcode,language=Python}
\lstdefinestyle{utpr}{style=utpcode,language=R,tabsize=2}
\lstdefinestyle{utpc}{style=utpcode,language=C}
\lstdefinestyle{utpcpp}{style=utpcode,language=C++}
\lstdefinestyle{utpbash}{style=utpcode,language=bash,
  rulecolor=\color{UTPListingTerminalFrame},aboveskip=0.6em,belowskip=0.7em}
\lstdefinestyle{utpcnumbered}{style=utpc,numbers=left,numberstyle=\tiny\color{UTPMuted},
  numbersep=7pt,xleftmargin=1.7em,framexleftmargin=1.4em}
\lstdefinestyle{utpbashnumbered}{style=utpbash,numbers=left,numberstyle=\tiny\color{UTPMuted},
  numbersep=7pt,xleftmargin=1.7em,framexleftmargin=1.4em}

% Pseudocode. A book that sets algorithms asks for them once, after the
% input, and gets algorithm2e loaded and coloured to match:
%
%     \utpusealgorithms
%
% \algostyle then sets one algorithm environment a size smaller than the
% text around it.
\newcommand{\utpalgorithmstyle}{%
  \SetAlCapNameFnt{\bfseries\color{UTPNavy}}%
  \SetAlCapFnt{\bfseries\color{UTPNavy}}%
  \SetAlgoCaptionSeparator{.}%
  \SetKwInput{KwIn}{Input}%
  \SetKwInput{KwOut}{Output}%
  \SetNlSty{color}{UTPMuted}{}%
  \SetCommentSty{\itshape\color{UTPMuted}}%
  \SetKwSty{bfseries\color{UTPBlue}}%
  \SetArgSty{}%
  \providecommand{\algostyle}{\small\SetAlgoNlRelativeSize{-1}\DontPrintSemicolon\SetAlgoVlined}%
}
\newcommand{\utpusealgorithms}{%
  \usepackage[linesnumbered,ruled,vlined]{algorithm2e}%
  \utpalgorithmstyle
}
